\documentclass[11pt, a4paper]{article}

% bfw-style.tex — gemeinsame Style-Definitionen für alle Handouts/Aufgabenhefte
% Voraussetzung: Kompilation mit xelatex (wegen fontspec)

\usepackage[ngerman]{babel}
% Babel-ngerman macht " zu einem aktiven Shorthand (z.B. für "a -> ä). In unseren
% Texten verwenden wir " als normales Zeichen — daher Shorthand komplett ausschalten.
\shorthandoff{"}
\usepackage{geometry}
\geometry{a4paper, top=2.4cm, bottom=2.4cm, left=2.3cm, right=2.3cm,
          headheight=15pt, headsep=0.7cm, footskip=1.1cm}

\usepackage{fontspec}
% Fallback-Kette: Source Sans Pro -> Calibri -> Segoe UI -> Latin Modern Sans
% (Latin Modern ist immer mit MiKTeX vorhanden)
\IfFontExistsTF{Source Sans Pro}{%
  \setmainfont{Source Sans Pro}[Ligatures=TeX]%
}{\IfFontExistsTF{Calibri}{%
  \setmainfont{Calibri}[Ligatures=TeX]%
}{\IfFontExistsTF{Segoe UI}{%
  \setmainfont{Segoe UI}[Ligatures=TeX]%
}{%
  \setmainfont{Latin Modern Sans}[Ligatures=TeX]%
}}}
\IfFontExistsTF{Source Code Pro}{%
  \setmonofont{Source Code Pro}[Scale=0.92]%
}{\IfFontExistsTF{Consolas}{%
  \setmonofont{Consolas}[Scale=0.92]%
}{%
  \setmonofont{Latin Modern Mono}[Scale=0.92]%
}}

% Gesperrte Versalien für Labels/Titelbalken (Kapitälchen-Ersatz, funktioniert
% mit jeder OpenType-Schrift — Latin Modern Sans hat keine echten Kapitälchen)
\newcommand{\bfwlabelfont}{\footnotesize\bfseries\addfontfeatures{LetterSpace=8.0}}

\usepackage{xcolor}
% Farbpalette: hell, freundlich, modern
\definecolor{bfwPrimary}{HTML}{0EA5A4}    % Petrol/Türkis - Primär
\definecolor{bfwPrimaryDark}{HTML}{0F766E} % dunkler Petrol für Headings
\definecolor{bfwAccent}{HTML}{F59E0B}     % Amber - Akzent
\definecolor{bfwSuccess}{HTML}{10B981}    % Grün - Erfolg / Lösung
\definecolor{bfwWarn}{HTML}{F97316}       % Orange - Achtung
\definecolor{bfwInk}{HTML}{0F172A}        % fast Schwarz
\definecolor{bfwInkSoft}{HTML}{475569}    % gedämpftes Grau
\definecolor{bfwBgSoft}{HTML}{F8FAFC}     % off-white
\definecolor{bfwBgTeal}{HTML}{ECFEFF}     % helles Türkis
\definecolor{bfwBgAmber}{HTML}{FEF3C7}    % helles Gelb
\definecolor{bfwBgRose}{HTML}{FFE4E6}     % helles Rosé für Eselsbrücken
\definecolor{bfwTabHead}{HTML}{E4F3F2}    % zarte Kopfzeilen-Hinterlegung für Tabellen

\usepackage[colorlinks=true, linkcolor=bfwPrimaryDark, urlcolor=bfwPrimaryDark]{hyperref}
\usepackage{lastpage}
\usepackage{enumitem}
\setlist[itemize]{leftmargin=*, itemsep=2pt, topsep=2pt}
\setlist[enumerate]{leftmargin=*, itemsep=2pt, topsep=2pt}

\usepackage[most]{tcolorbox}
\usepackage{booktabs}
\usepackage{colortbl}
\usepackage{tikz}
\usetikzlibrary{decorations.pathmorphing, positioning, shapes.geometric, arrows.meta}

% ---------------------------------------------------------------------------
% Einheitliches Tabellen-Design
%   - etwas mehr Zeilenluft, dezente graue Linien (wirkt auch mit \hline ruhiger)
%   - Kopfzeile einer Tabelle bei Bedarf zart hinterlegen: \rowcolor{bfwTabHead}
% ---------------------------------------------------------------------------
\renewcommand{\arraystretch}{1.25}
\arrayrulecolor{bfwInkSoft!50}
\setlength{\heavyrulewidth}{1pt}
\setlength{\lightrulewidth}{0.5pt}

% ---------------------------------------------------------------------------
% Boxen — klare farbige Titelbalken mit gesperrten Versal-Labels
% (Titel bleibt per Option überschreibbar: \begin{merksatz}[title={...}])
% ---------------------------------------------------------------------------
\tcbset{bfwbox/.style={%
  enhanced, breakable,
  boxrule=0.5pt, arc=2.5pt,
  left=10pt, right=10pt, top=8pt, bottom=8pt,
  toptitle=4pt, bottomtitle=4pt,
  titlerule=0pt,
  fonttitle=\bfwlabelfont,
  coltitle=white
}}

% Merkkästchen — wichtige Definition / Kernpunkt
\newtcolorbox{merksatz}[1][]{%
  bfwbox,
  colback=bfwBgTeal,
  colframe=bfwPrimaryDark,
  colbacktitle=bfwPrimaryDark,
  title={MERKE},
  #1
}

% Eselsbrücke — Mnemoniks
\newtcolorbox{eselsbruecke}[1][]{%
  bfwbox,
  colback=bfwBgRose,
  colframe=bfwAccent!85!black,
  colbacktitle=bfwAccent!85!black,
  title={ESELSBRÜCKE},
  #1
}

% Beispiel-Box
\newtcolorbox{beispiel}[1][]{%
  bfwbox,
  colback=bfwBgSoft,
  colframe=bfwInkSoft,
  colbacktitle=bfwInkSoft,
  title={BEISPIEL},
  #1
}

% Lösungs-Box (für Aufgabenhefte)
\newtcolorbox{loesung}[1][]{%
  bfwbox,
  colback=bfwSuccess!7,
  colframe=bfwSuccess!65!black,
  colbacktitle=bfwSuccess!65!black,
  title={LÖSUNG},
  #1
}

% Achtung-Box — typische Fehler / Stolperfallen
\newtcolorbox{achtung}[1][]{%
  bfwbox,
  colback=bfwWarn!6,
  colframe=bfwWarn!85!black,
  colbacktitle=bfwWarn!85!black,
  title={ACHTUNG},
  #1
}

% Formel-Box — für zentrale Formeln (groß, ruhig, ohne Titelbalken)
\newtcolorbox{formelbox}[1][]{%
  enhanced,
  colback=bfwBgTeal,
  colframe=bfwPrimaryDark,
  boxrule=1pt, arc=3pt,
  left=10pt, right=10pt, top=6pt, bottom=6pt,
  #1
}

% Glossar-Eintrag
\newcommand{\glossareintrag}[2]{%
  \par\noindent\textbf{\color{bfwPrimaryDark}#1} \enspace #2\par\smallskip
}

% ---------------------------------------------------------------------------
% Abschnitts-Design: farbiges Nummern-Quadrat + feine Linie unter \section,
% dezent farbige \subsection/\subsubsection
% ---------------------------------------------------------------------------
\usepackage{titlesec}
\newcommand{\bfwSecBadge}[1]{%
  \begin{tikzpicture}[baseline={([yshift=-0.62em]n.center)}]
    \node[fill=bfwPrimaryDark, text=white, font=\normalsize\bfseries,
          minimum width=1.55em, minimum height=1.55em, inner sep=1pt,
          rounded corners=1.5pt] (n) {#1};
  \end{tikzpicture}}
\titleformat{\section}
  {\Large\bfseries\color{bfwPrimaryDark}}
  {\bfwSecBadge{\thesection}}{0.6em}{}
  [\vspace{-0.2em}{\color{bfwPrimary!60}\titlerule[0.8pt]}]
\titleformat{\subsection}
  {\large\bfseries\color{bfwPrimaryDark}}
  {\textcolor{bfwPrimary}{\thesubsection}}{0.5em}{}
\titleformat{\subsubsection}
  {\normalsize\bfseries\color{bfwInk}}
  {\textcolor{bfwPrimary}{\thesubsubsection}}{0.4em}{}
\titlespacing*{\section}{0pt}{1.6em}{1em}
\titlespacing*{\subsection}{0pt}{1.2em}{0.5em}

% TikZ-Stil "skizzenhaft" — etwas weicher als Rough.js, aber stabil
\tikzset{
  roughbox/.style = {
    draw=bfwPrimaryDark, thick, fill=bfwBgTeal,
    rounded corners=4pt, inner sep=6pt
  },
  roughaccent/.style = {
    draw=bfwAccent, thick, fill=bfwBgAmber,
    rounded corners=4pt, inner sep=6pt
  },
  roughline/.style = {
    draw=bfwInkSoft, thick, -{Stealth[length=2.5mm]}
  }
}

% Bullet-Symbole (bewusst kein FontAwesome — vermeidet Font-Installations-Probleme)
\newcommand{\faLightbulb}{$\bullet$}
\newcommand{\faBookmark}{$\bullet$}

% ---------------------------------------------------------------------------
% Kopf-/Fußzeile
%   Kopf links:  Wortlogo „Wirtschaftsfachwirt"
%   Kopf rechts: Dokument-Kurztext, gesetzt via \kopfzeile{...}
%   Fuß links:   © Can Siebert · 2026 — Fuß rechts: Seite X von Y
%   Titelseite (Seite 1) bleibt ohne Kopf-/Fußzeile (\handouttitel setzt
%   \thispagestyle{empty}).
% ---------------------------------------------------------------------------
\usepackage{fancyhdr}
\newcommand{\bfwKopfzeileText}{}
\newcommand{\kopfzeile}[1]{\renewcommand{\bfwKopfzeileText}{#1}}
\pagestyle{fancy}
\fancyhf{}
\fancyhead[L]{\small\bfseries\color{bfwPrimaryDark}Wirtschaftsfachwirt}
\fancyhead[R]{\small\color{bfwInkSoft}\bfwKopfzeileText}
\renewcommand{\headrulewidth}{0.6pt}
\renewcommand{\headrule}{{\color{bfwPrimary}\hrule width\headwidth height 0.6pt}}
\fancyfoot[L]{\small\color{bfwInkSoft}\copyright\ Can Siebert\,\textperiodcentered\,2026}
\fancyfoot[R]{\small\color{bfwInkSoft}Seite \thepage\ von \pageref*{LastPage}}
% Auch \thispagestyle{plain} (z.B. durch \maketitle) soll leer bleiben:
\fancypagestyle{plain}{\fancyhf{}\renewcommand{\headrulewidth}{0pt}\renewcommand{\headrule}{}}

% ---------------------------------------------------------------------------
% \handouttitel{Obertitel}{Untertitel}{Kurs-Label}{Termin-Zeile}
%   Einheitlicher professioneller Titelblock: Dachzeile, farbige Headline,
%   Untertitel, farbige Trennlinie und dezente Meta-Zeile (Kurs/Termin/Dozent).
%   Ersetzt \title/\maketitle. Direkt nach \begin{document} aufrufen.
% ---------------------------------------------------------------------------
\newcommand{\handouttitel}[4]{%
  \thispagestyle{empty}%
  \begingroup
  \setlength{\parindent}{0pt}%
  {\bfwlabelfont\color{bfwPrimary}WIRTSCHAFTSFACHWIRT (IHK)\par}
  \vspace{0.55em}
  {\huge\bfseries\color{bfwPrimaryDark}#1\par}
  \vspace{0.5em}
  {\large\color{bfwInkSoft}#2\par}
  \vspace{0.9em}
  {\color{bfwPrimary}\rule{\linewidth}{2pt}}\par
  \vspace{0.55em}
  {\small\color{bfwInkSoft}#3\ \,\textperiodcentered\,\ #4\ \,\textperiodcentered\,\ Dozent: Can Siebert\par}
  \endgroup
  \vspace{1.4em}%
}


\kopfzeile{Aufgabenheft BWL Lektion 3}

\begin{document}
\handouttitel{Aufgabenheft Lektion~3 (BWL)}%
  {Existenzgründung \& Unternehmensrechtsformen --- Übungen im IHK-Klausurformat}%
  {1.3 Betriebswirtschaftliche Grundlagen}%
  {Selbstlernmaterial zur Lektion vom Sa, 26.09.2026}


\begin{merksatz}[title={\bfseries So nutzen Sie dieses Aufgabenheft}]
Bearbeiten Sie alle Aufgaben zunächst \emph{ohne} in die Lösungen zu schauen. Schreiben
Sie Ihre Antworten in vollständigen Sätzen --- die IHK-Prüfung verlangt formulierte
Begründungen, nicht bloße Stichworte. Beachten Sie die \emph{Operatoren} (nennen,
erläutern, beurteilen, begründen, empfehlen) --- sie geben den geforderten
Antwort-Tiefegrad vor. Erst danach öffnen Sie den Lösungsteil ab
Seite~\pageref{loesungen} und vergleichen. Ergänzend: Übungsband Betriebswirtschaft,
Kapitel~2 (Übungen 45--53). Wegen der langen Pause bis Lektion~4 (27.02.2027) lohnt es
sich, dieses Heft nach einigen Wochen ein zweites Mal zu bearbeiten.
\end{merksatz}

\tableofcontents
\vspace{1em}
\hrule

\section{Aufgaben}

\subsection*{Aufgabe~1 --- Gründungsphasen (8 Punkte)}

\begin{enumerate}[label=(\alph*)]
  \item \textbf{Bringen} Sie die folgenden Phasen einer Unternehmensgründung \textbf{in
    die richtige Reihenfolge}: Orientierung --- Festigung --- Geschäftsidee ---
    Konzeption --- Umsetzung. \hfill (3~P.)
  \item \textbf{Beschreiben} Sie kurz, was in der Konzeptionsphase und was in der
    Festigungsphase geschieht. \hfill (4~P.)
  \item \textbf{Nennen} Sie einen Grund, warum die Zahl der Existenzgründungen in Jahren
    mit guter Arbeitsmarktlage tendenziell sinkt. \hfill (1~P.)
\end{enumerate}

\subsection*{Aufgabe~2 --- Voraussetzungen der Existenzgründung (12 Punkte)}

Frau Yilmaz, Konditormeisterin mit 12 Jahren Berufserfahrung, möchte sich mit einer
eigenen Confiserie selbstständig machen.

\begin{enumerate}[label=(\alph*)]
  \item \textbf{Nennen} Sie je drei persönliche Voraussetzungen, drei fachliche
    Kompetenzen und drei unternehmerische (kaufmännische) Kompetenzen, die Frau Yilmaz
    vor der Gründung prüfen sollte. \hfill (9~P.)
  \item \textbf{Erläutern} Sie, warum die finanzielle Basis vor der Gründung realistisch
    geprüft werden muss. \hfill (3~P.)
\end{enumerate}

\subsection*{Aufgabe~3 --- Businessplan (14 Punkte)}

Frau Yilmaz benötigt für Ladenausbau, Maschinen und Anlaufphase rund 80.000~€
Startkapital und möchte ihre Hausbank von der Finanzierung überzeugen.

\begin{enumerate}[label=(\alph*)]
  \item \textbf{Erklären} Sie den Begriff Businessplan und \textbf{erläutern} Sie seine
    drei Funktionen (mit Adressaten). \hfill (5~P.)
  \item \textbf{Skizzieren} Sie den Aufbau des Businessplans von Frau Yilmaz mit sechs
    inhaltlichen Bestandteilen und je einem Stichwort zum Inhalt. \hfill (6~P.)
  \item \textbf{Nennen} Sie drei Teilpläne der Finanzplanung. \hfill (3~P.)
\end{enumerate}

\subsection*{Aufgabe~4 --- Fördermöglichkeiten (6 Punkte)}

\begin{enumerate}[label=(\alph*)]
  \item \textbf{Erläutern} Sie, wie Existenzgründer über die KfW gefördert werden können
    (einschließlich Hausbankprinzip). \hfill (3~P.)
  \item \textbf{Erläutern} Sie den Gründungszuschuss: Träger, Zielgruppe und eine
    wesentliche Voraussetzung. \hfill (3~P.)
\end{enumerate}

\subsection*{Aufgabe~5 --- Personengesellschaften: richtig oder falsch? (10 Punkte)}

\textbf{Beurteilen} Sie die folgenden Aussagen über Personengesellschaften und
\textbf{stellen} Sie falsche Aussagen \textbf{richtig}:

\begin{enumerate}[label=(\alph*)]
  \item Die Eigentümer haften unbeschränkt mit ihrem Betriebs- und Privatvermögen. \hfill (2~P.)
  \item Bei Personengesellschaften handelt es sich um juristische Personen. \hfill (2~P.)
  \item Die Gesellschaft selbst unterliegt der Einkommensteuerpflicht. \hfill (2~P.)
  \item Die Kommanditgesellschaft ist eine Personengesellschaft. \hfill (2~P.)
  \item Eine GbR entsteht nur durch notariellen Vertrag und Eintragung ins
    Handelsregister. \hfill (2~P.)
\end{enumerate}

\subsection*{Aufgabe~6 --- Kommanditgesellschaft (10 Punkte)}

\begin{enumerate}[label=(\alph*)]
  \item \textbf{Unterscheiden} Sie Komplementär und Kommanditist nach den Kriterien
    Haftung, Geschäftsführung und Vertretung. \hfill (6~P.)
  \item \textbf{Erläutern} Sie zwei Rechte, die dem Kommanditisten trotz seines
    Ausschlusses von der Geschäftsführung zustehen. \hfill (2~P.)
  \item \textbf{Begründen} Sie, warum sich die KG eignet, um Kapitalgeber zu gewinnen,
    ohne Führungsrechte abzugeben. \hfill (2~P.)
\end{enumerate}

\subsection*{Aufgabe~7 --- GmbH und UG (haftungsbeschränkt) (12 Punkte)}

\begin{enumerate}[label=(\alph*)]
  \item \textbf{Vervollständigen} Sie den Lückentext (8 Lücken):\\[0.3em]
    Die GmbH ist eine \rule{3cm}{0.4pt}. Für die Verbindlichkeiten steht als Deckung
    lediglich das \rule{3cm}{0.4pt}, d.\,h.\ das \rule{3cm}{0.4pt}, zur Verfügung. Es
    beträgt mindestens \rule{2cm}{0.4pt}~€. Es können auch \rule{3cm}{0.4pt} bei der
    Gründung eingebracht werden. Die Geschäftsführung steht den Gesellschaftern offen,
    kann aber auch durch einen externen \rule{3cm}{0.4pt} erfolgen. Die Gewinn- und
    Verlustverteilung erfolgt im Verhältnis der Einlagen, wenn nichts anderes im
    \rule{3cm}{0.4pt} geregelt ist. Das Kontrollorgan der GmbH ist die
    \rule{3.6cm}{0.4pt}. \hfill (8~P.)
  \item \textbf{Nennen} Sie vier Besonderheiten der UG (haftungsbeschränkt) im Vergleich
    zur GmbH. \hfill (4~P.)
\end{enumerate}

\subsection*{Aufgabe~8 --- Organe der AG (9 Punkte)}

\textbf{Stellen} Sie für die drei Organe der Aktiengesellschaft --- Vorstand,
Aufsichtsrat, Hauptversammlung --- jeweils die Zusammensetzung und die Aufgabe
\textbf{dar}. \hfill (je 3~P.)

\subsection*{Aufgabe~9 --- Fallaufgabe: Welche Rechtsform passt? (19 Punkte)}

Drei Gründungsfälle:

\begin{itemize}
  \item \textbf{Fall A:} Zwei Ingenieurinnen wollen gemeinsam ein Unternehmen für die
    Entwicklung von Maschinenbau-Prototypen gründen. Das Geschäft ist technisch riskant
    (Gewährleistung, Produkthaftung). Beide bringen je 15.000~€ mit, wollen beide die
    Geschäfte führen und ihr Privatvermögen schützen.
  \item \textbf{Fall B:} Ein Handwerksmeister will seinen gut laufenden Betrieb
    vergrößern. Sein Onkel stellt 100.000~€ zur Verfügung, möchte aber weder mitarbeiten
    noch unbeschränkt haften. Der Meister will die Geschäfte weiterhin allein führen und
    ist bereit, persönlich zu haften.
  \item \textbf{Fall C:} Ein Familienunternehmen (Großhandel) sucht eine Rechtsform, die
    die organisatorischen und steuerlichen Vorteile einer Personengesellschaft bietet,
    bei der aber keine natürliche Person unbeschränkt haftet.
\end{itemize}

\begin{enumerate}[label=(\alph*)]
  \item \textbf{Empfehlen} Sie für jeden Fall eine passende Rechtsform und
    \textbf{begründen} Sie Ihre Empfehlung jeweils anhand von mindestens drei
    Vergleichskriterien (z.\,B. Haftung, Geschäftsführung, Kapital, Gewinnverteilung).
    \hfill (12~P.)
  \item \textbf{Erläutern} Sie die rechtliche Konstruktion der in Fall~C empfohlenen
    Rechtsform. \hfill (4~P.)
  \item \textbf{Prüfen} Sie anhand der Firmengrundsätze, ob die Firma „Prototec
    Maschinenbau GmbH`` für Fall~A zulässig wäre. \hfill (3~P.)
\end{enumerate}

\newpage
\section{Lösungen}\label{loesungen}

\begin{loesung}[title={Lösung Aufgabe~1 --- Gründungsphasen}]
\textbf{(a)} 1.~Geschäftsidee --- 2.~Orientierung --- 3.~Konzeption --- 4.~Umsetzung ---
5.~Festigung.

\textbf{(b)} \emph{Konzeptionsphase:} Aus der Idee wird ein tragfähiges Konzept --- der
Businessplan entsteht (Marktanalyse, Standort, Marketing, Rechtsformwahl, Kapitalbedarfs-
und Finanzierungsplanung); Gespräche mit Kapitalgebern werden vorbereitet.
\emph{Festigungsphase:} Das junge Unternehmen konsolidiert sich am Markt --- Ausbau des
Kundenstamms, Sicherung der Liquidität, Professionalisierung der Abläufe, laufender
Soll-Ist-Vergleich mit dem Businessplan und ggf.\ Nachsteuern.

\textbf{(c)} Bei guter Arbeitsmarktlage finden potenzielle Gründer leichter eine sichere
Festanstellung und verzichten auf das Gründungsrisiko (verstärkt durch Fachkräftemangel
und demografische Entwicklung).
\end{loesung}

\begin{loesung}[title={Lösung Aufgabe~2 --- Voraussetzungen}]
\textbf{(a)} \emph{Persönlich:} sich selbst Ziele setzen und konsequent verfolgen; Mut
zum Risiko und Umgang mit Rückschlägen; Rückhalt der Familie und Bereitschaft zu hohem
zeitlichen Einsatz; Disziplin im Umgang mit Geld (drei Nennungen genügen).
\emph{Fachlich:} Ausbildung/Berufserfahrung passen zum Vorhaben (Konditormeisterin:
gegeben); Branchenkenntnis (Trends, Kundschaft, Wettbewerb); Bereitschaft, Defizite
auszugleichen; Umgang mit branchenspezifischer EDV (Kassensystem, Warenwirtschaft).
\emph{Unternehmerisch/kaufmännisch:} Kenntnisse in Kalkulation, Buchführung, Steuern und
Recht; selbstbewusstes, kontaktfreudiges Auftreten; Fähigkeit zu delegieren; Kreativität
für die Weiterentwicklung des Unternehmens.

\textbf{(b)} Finanzierungsmängel sind eine Hauptursache des Scheiterns junger
Unternehmen: Wird der Kapitalbedarf (Investitionen, Anlaufverluste, private
Lebenshaltung) unterschätzt oder fehlt Eigenkapital (Faustregel 15--20\,\%), droht schon
in der Startphase Zahlungsunfähigkeit --- unabhängig von der Qualität der Geschäftsidee.
\end{loesung}

\begin{loesung}[title={Lösung Aufgabe~3 --- Businessplan}]
\textbf{(a)} Der Businessplan ist die \emph{schriftliche} Ausarbeitung zur Darlegung der
Geschäftsidee, ihrer Realisierung und der daraus resultierenden Ergebnisse. Funktionen:
(1)~Er richtet sich v.\,a.\ an die \emph{Geldgeber} (hier: die Hausbank) und künftige
Geschäftspartner --- Entscheidungsgrundlage für die Finanzierung. (2)~Er ist ein Mittel
für die Existenzgründerin, ihre Geschäftsidee \emph{planmäßig} zu durchdenken.
(3)~Er dient als \emph{Führungsgrundlage} für Mitarbeiter und Gremien im zu gründenden
Unternehmen (später: Soll-Ist-Kontrolle).

\textbf{(b)} (1)~\emph{Executive Summary} --- Kurzfassung des Vorhabens;
(2)~\emph{Geschäftsmodell/Geschäftsidee} --- Confiserie mit handgefertigten Pralinen,
Kundennutzen, Gründerprofil; (3)~\emph{Markt und Wettbewerb} --- Zielgruppen, Standort,
Konkurrenzanalyse; (4)~\emph{Marketing} --- Sortiment, Preise, Vertriebswege, Werbung;
(5)~\emph{Organisation und Rechtsform} --- Personal, gewählte Rechtsform mit Begründung;
(6)~\emph{Finanzplanung} --- Kapitalbedarf 80.000~€, Finanzierung, Liquidität,
Ertragsvorschau; ergänzend: Chancen und Risiken.

\textbf{(c)} Z.\,B. Kapitalbedarfsplan, Investitionsplan, Finanzierungsplan,
Liquiditätsplan, Umsatzschätzung/Ertragsvorschau (Rentabilitätsplan) --- drei Nennungen
genügen.
\end{loesung}

\begin{loesung}[title={Lösung Aufgabe~4 --- Fördermöglichkeiten}]
\textbf{(a)} Die staatliche KfW-Bank vergibt zinsgünstige Förderkredite für
Existenzgründer (z.\,B. ERP-Gründerkredit „StartGeld`` bis 125.000~€). Der Antrag läuft
nach dem \emph{Hausbankprinzip} über die eigene Geschäftsbank; die KfW refinanziert den
Kredit und übernimmt einen großen Teil des Ausfallrisikos (Haftungsfreistellung), was die
Kreditvergabe an Gründer erleichtert.

\textbf{(b)} Der \emph{Gründungszuschuss} wird von der \emph{Agentur für Arbeit} als
Ermessensleistung an Empfänger von \emph{Arbeitslosengeld~I} gezahlt, die sich
hauptberuflich selbstständig machen. Wesentliche Voraussetzung ist u.\,a.\ der Nachweis
der Tragfähigkeit der Gründung durch eine fachkundige Stelle (z.\,B. IHK); gezahlt wird
der bisherige ALG-I-Betrag plus eine Pauschale zur sozialen Absicherung.
\end{loesung}

\begin{loesung}[title={Lösung Aufgabe~5 --- Personengesellschaften}]
\textbf{(a) Richtig.} Kennzeichen der Personengesellschaft ist die unbeschränkte,
persönliche Haftung (bei der KG allerdings nur für den Komplementär).

\textbf{(b) Falsch.} Personengesellschaften sind \emph{keine} juristischen Personen ---
Träger von Rechten und Pflichten sind (auch) die Gesellschafter persönlich.

\textbf{(c) Falsch.} Nicht die Gesellschaft, sondern die \emph{Gesellschafter} versteuern
ihren Gewinnanteil mit ihrer persönlichen Einkommensteuer.

\textbf{(d) Richtig.} Die KG ist eine Personengesellschaft mit zwei Gesellschaftergruppen
(Komplementäre und Kommanditisten).

\textbf{(e) Falsch.} Die GbR entsteht durch \emph{formfreien} Gesellschaftsvertrag
(auch mündlich); ein Handelsregistereintrag ist nicht erforderlich (seit 2024 ist eine
freiwillige Eintragung als eGbR ins Gesellschaftsregister möglich).
\end{loesung}

\begin{loesung}[title={Lösung Aufgabe~6 --- Kommanditgesellschaft}]
\textbf{(a)} \emph{Komplementär (Vollhafter):} haftet unbeschränkt, unmittelbar und
gesamtschuldnerisch mit seinem gesamten Vermögen; ihm obliegen Geschäftsführung und
Vertretung der KG. \emph{Kommanditist (Teilhafter):} haftet nur bis zur Höhe seiner im
Handelsregister eingetragenen Haftsumme (Einlage); von Geschäftsführung und Vertretung
ist er ausgeschlossen.

\textbf{(b)} Der Kommanditist hat \emph{Kontrollrechte} (Einsicht in den Jahresabschluss
und die Bücher) und ein \emph{Widerspruchsrecht} bei außergewöhnlichen Geschäften, die
über den gewöhnlichen Betrieb hinausgehen.

\textbf{(c)} Kapitalgeber können als Kommanditisten mit begrenztem Risiko (nur die
Einlage) beteiligt werden, ohne dass sie Geschäftsführungs- oder Vertretungsbefugnisse
erhalten --- die Leitung bleibt vollständig beim Komplementär.
\end{loesung}

\begin{loesung}[title={Lösung Aufgabe~7 --- GmbH und UG}]
\textbf{(a)} Die GmbH ist eine \emph{Kapitalgesellschaft}. Für die Verbindlichkeiten
steht als Deckung lediglich das \emph{Gesellschaftsvermögen}, d.\,h.\ das
\emph{Stammkapital}, zur Verfügung. Es beträgt mindestens \emph{25.000}~€. Es können auch
\emph{Sacheinlagen} bei der Gründung eingebracht werden. Die Geschäftsführung steht den
Gesellschaftern offen, kann aber auch durch einen externen \emph{Geschäftsführer}
erfolgen. Die Gewinn- und Verlustverteilung erfolgt im Verhältnis der Einlagen, wenn
nichts anderes im \emph{Gesellschaftsvertrag} geregelt ist. Das Kontrollorgan der GmbH
ist die \emph{Gesellschafterversammlung}. (Ergänzung: Ab 500 Mitarbeitern existiert
zusätzlich ein Aufsichtsrat.)

\textbf{(b)} (1)~Gründung ab 1~€ Stammkapital möglich; (2)~Sacheinlagen bei der Gründung
unzulässig (Bareinzahlung in voller Höhe); (3)~Pflicht zur gesetzlichen Rücklage: 25\,\%
des Jahresüberschusses, bis 25.000~€ erreicht sind (danach Umfirmierung zur GmbH
möglich); (4)~Rechtsformzusatz „UG (haftungsbeschränkt)`` muss ungekürzt geführt werden.
\end{loesung}

\begin{loesung}[title={Lösung Aufgabe~8 --- Organe der AG}]
\emph{Vorstand:} eine oder mehrere Personen --- leitendes Organ zur Wahrnehmung der
Geschäftsführung; leitet die Gesellschaft in eigener Verantwortung und vertritt sie nach
außen.
\emph{Aufsichtsrat:} mindestens drei Mitglieder --- Kontrollorgan; bestellt, überwacht
und berät den Vorstand und prüft den Jahresabschluss.
\emph{Hauptversammlung:} die Aktionäre --- beschließendes Organ; entscheidet u.\,a.\ über
Gewinnverwendung (Dividende), Entlastung von Vorstand und Aufsichtsrat, Wahl der
Anteilseignervertreter im Aufsichtsrat und Satzungsänderungen.
\end{loesung}

\begin{loesung}[title={Lösung Aufgabe~9 --- Fallaufgabe Rechtsformwahl}]
\textbf{(a) Fall A: GmbH.} \emph{Haftung:} Das riskante Prototypengeschäft spricht für
die Beschränkung auf das Gesellschaftsvermögen --- das Privatvermögen bleibt geschützt.
\emph{Kapital:} 2~$\times$~15.000~€ = 30.000~€ übersteigen das Mindeststammkapital von
25.000~€. \emph{Geschäftsführung:} Beide Gründerinnen können zu Geschäftsführerinnen
bestellt werden. (Alternative bei knapperem Kapital: UG (haftungsbeschränkt); eine GbR/OHG
scheidet wegen der unbeschränkten persönlichen Haftung aus.)

\textbf{Fall B: KG.} \emph{Haftung/Leitung:} Der Meister wird Komplementär --- er haftet
persönlich (dazu ist er bereit) und behält Geschäftsführung und Vertretung allein.
\emph{Kapital:} Der Onkel wird Kommanditist: Er bringt 100.000~€ ein, haftet nur mit
dieser Einlage und ist von der Geschäftsführung ausgeschlossen --- genau sein Wunsch,
weder mitzuarbeiten noch unbeschränkt zu haften. \emph{Gewinnverteilung:} flexibel im
Gesellschaftsvertrag regelbar.

\textbf{Fall C: GmbH \& Co. KG.} \emph{Haftung:} Keine natürliche Person haftet
unbeschränkt. \emph{Rechtsnatur:} Es bleibt eine Personengesellschaft (KG) mit deren
organisatorischen und steuerlichen Eigenschaften (Einkommensbesteuerung auf
Gesellschafterebene, flexible Vertragsgestaltung). \emph{Leitung:} Die Familie steuert
das Unternehmen über die Geschäftsführung der Komplementär-GmbH.

\textbf{(b)} Bei der GmbH \& Co. KG ist der einzige \emph{Komplementär} (Vollhafter)
keine natürliche Person, sondern eine \emph{GmbH}. Diese haftet zwar unbeschränkt für
die KG --- aber ihrerseits nur mit ihrem Gesellschaftsvermögen. Die Familienmitglieder
sind \emph{Kommanditisten} (Haftung nur mit der Einlage) und meist zugleich
Gesellschafter der Komplementär-GmbH; die Geschäfte führt der Geschäftsführer der GmbH.

\textbf{(c)} Die Firma ist \emph{unzulässig}: Der Rechtsformzusatz „GmbH`` verstößt gegen
den Grundsatz der \emph{Firmenwahrheit}, solange tatsächlich keine GmbH, sondern z.\,B.
eine GbR oder OHG vorliegt. Zulässig wäre der Zusatz erst nach notarieller Gründung und
\emph{konstitutiver} Eintragung der GmbH ins Handelsregister~B. Der Sach-/Fantasiename
„Prototec Maschinenbau`` selbst wäre unbedenklich (unterscheidungskräftig, nicht
irreführend).
\end{loesung}

\vspace{1em}
\hrule
\vspace{0.8em}
\noindent\textsf{\small Weiterüben: Übungsband Betriebswirtschaft Kap.~2 (mit Lösungsteil) und das
interaktive Quiz dieser Lektion. Nächste Lektion (Sa, 27.02.2027): Unternehmenszusammenschlüsse.}

\end{document}
